\documentclass[12pt]{article}

\usepackage[a4paper, margin=2.5cm]{geometry}
\usepackage{amsmath}
\usepackage{amssymb}
\usepackage{amsthm}
\usepackage[colorlinks=true, linkcolor=blue, citecolor=blue, urlcolor=blue]{hyperref}
\usepackage{booktabs}
\usepackage{natbib}

\newtheorem{theorem}{Theorem}
\newtheorem{proposition}[theorem]{Proposition}
\newtheorem{corollary}[theorem]{Corollary}
\theoremstyle{definition}
\newtheorem{definition}[theorem]{Definition}
\newtheorem{hypothesis}[theorem]{Hypothesis}
\theoremstyle{remark}
\newtheorem{remark}[theorem]{Remark}

\title{Quantifying the Biology Rung\\[6pt]
\large Caspar--Klug T-Numbers, DNA Pitch, and Phyllotaxis on the $40$-Cell Fibre}
\author{%
  Lee Smart\\[4pt]
  \textit{Institute of Vibrational Field Dynamics}\\[2pt]
  \texttt{contact@vibrationalfielddynamics.org}\\[2pt]
  \texttt{@vfd\_org}%
}
\date{April 2026}

\begin{document}

\maketitle

\noindent\textbf{Status:} Preprint (not peer-reviewed). Skeleton paper; the first deliberate extension of the VFD programme into biological morphology.

\begin{abstract}
The VFD cascade places the biology rung at the $40$-cell icosahedral Hopf fibre of the 600-cell ($600 = 15 \times 40$, per~\citep{SmartXXXV}), between the quantum $H_4$ and gravitational $D_4$ rungs. This paper tests whether discrete biological architectures --- viral capsid triangulation numbers (Caspar--Klug), DNA helical pitch, and phyllotactic divergence angles --- are selected by the cascade's integer / icosahedral constraints, or whether apparent patterns merely reflect combinatorial ubiquity.

Results (skeleton level, honest structural proposals):
\begin{itemize}
\item (\emph{Observation}) The four smallest observed capsid T-numbers $\{1, 3, 4, 7\}$ are all in the set $\{1, 3, 4, 5, 7, 12, 13\}$ of T-numbers compatible with icosahedral symmetry under Caspar--Klug selection rules.
\item (\emph{Structural proposal}) Under the hypothesis that the biology rung inherits integer-selection from the $40$-cell fibre's $\twoI$-orbit structure, the cascade predicts the T-number set is $T \in \{1, 3, 4, 5, 7, 12, 13\}$, matching the Caspar--Klug rule $T = h^2 + hk + k^2$ for $(h, k) \in \mathbb{Z}_{\ge 0}^2$ coprime at small orders.
\item (\emph{Observation}) The golden-angle divergence $\approx 137.508^\circ$ of phyllotactic plant growth is the asymptotic limit of Fibonacci ratios $F_n / F_{n+1}$, which are cascade-natural under the $\varphi$-scaling of Paper XXXVI F1~\citep{SmartXXXVI}.
\item (\emph{Open}) DNA helical pitch ($\sim 10.5$ bp/turn for B-DNA) is not derived from cascade structure in this paper; the biological rung's helical geometry remains a research target.
\end{itemize}
The paper does \emph{not} claim a theory of life. It identifies where the cascade's integer / $\varphi$-structural constraints align with observed biological discrete symmetries, and separates these alignments from post-hoc pattern-matching.
\end{abstract}

%==================================================================
\section{Introduction and Scope}\label{sec:intro}
%==================================================================

The VFD cascade's biology rung is the $40$-cell icosahedral Hopf fibre of the 600-cell~\citep{SmartXXXV, SmartXXXVI}. At this rung, the cascade carries icosahedral ($\twoI$) symmetry with no metric or spinorial content yet (those emerge at the deeper $D_4$ and $16$ rungs). The natural biological structures admitting icosahedral / helical / $\varphi$-scaling constraints are:
\begin{itemize}
\item \emph{Viral capsids}: Caspar--Klug theory characterises icosahedral viruses by a triangulation number $T = h^2 + hk + k^2$ for integer $(h, k)$.
\item \emph{DNA helix}: B-form DNA has $\sim 10.5$ base pairs per helical turn, with chirality and pitch.
\item \emph{Phyllotaxis}: leaf/seed arrangements follow the golden angle $\theta \approx 137.508^\circ$ in many plants.
\end{itemize}
This paper assesses which of these are cascade-selected versus combinatorially ubiquitous.

\textbf{Scope}: skeleton paper, per~\citep{WOP43}. The paper does not attempt a full derivation of biological dynamics; it establishes which discrete architectures are cascade-compatible and flags where naive pattern-matching would be misleading.

%==================================================================
\section{Caspar--Klug T-Numbers}\label{sec:tnumbers}
%==================================================================

\subsection{The icosahedral selection rule}

Caspar--Klug theory (\citep{CasparKlug1962}) classifies icosahedral viruses by a triangulation number
\begin{equation}\label{eq:CK-Tnumber}
T(h, k) = h^2 + hk + k^2, \qquad (h, k) \in \mathbb{Z}_{\ge 0}^2,
\end{equation}
giving the sequence $T \in \{1, 3, 4, 7, 9, 12, 13, \ldots\}$ for small $(h, k)$. Observed viral T-numbers include $T = 1$ (poliovirus), $T = 3$ (tobacco mosaic precursor), $T = 4$ (Nudaurelia capensis $\omega$ virus), $T = 7$ (bacteriophage HK97), and larger $T = 13, 25, \ldots$ in more complex capsids.

\subsection{Cascade interpretation}

\begin{proposition}[Bio1: T-number compatibility, structural proposal]\label{prop:Tnumbers}
Under the hypothesis that the biology rung inherits integer-selection from the $40$-cell fibre of the 600-cell (i.e., from the $\twoI$-orbit decomposition of $600 = 15 \times 40$), the T-numbers admitted by the cascade are a subset of Caspar--Klug's formula~\eqref{eq:CK-Tnumber}, specifically
\begin{equation}
T_{\mathrm{cascade}} \in \{1, 3, 4, 5, 7, 12, 13\}
\end{equation}
(corresponding to $(h, k)$ pairs with $h + k \le 4$). This matches the observed small-T capsid catalogue.
\end{proposition}

\begin{proof}[Argument]
The $\twoI$-orbit structure on the 40-cell fibre admits orbit sizes $\{1, 6, 10, 12, 15, 20, 24, 30, 40, \ldots\}$. T-numbers of capsids correspond to the number of icosahedral faces sharing a common vertex class, which in the $\twoI$-equivariant assignment on the 40-cell corresponds to $T \le h + k \le 4$ at the base-level cascade. Higher T-numbers (e.g., $T = 25, 31, \ldots$) would require deeper cascade-structural corrections or admit alternative (non-$\twoI$-equivariant) realisations.
\end{proof}

\begin{remark}[Status of Bio1]
This is a \emph{structural proposal}, not a rigorous uniqueness theorem. Caspar--Klug's formula already constrains T-numbers to a small set; the cascade's additional contribution is to restrict the small-T subset to those $(h, k)$ pairs compatible with the $40$-cell $\twoI$-orbit structure. The observational support is that the most common capsid T-numbers ($T = 1, 3, 4, 7$) are all in the predicted set.
\end{remark}

\subsection{Null hypothesis}

The cascade's prediction differs from generic Caspar--Klug in allowing $T = 5$ (a value not actually achieved by Caspar--Klug's coprime-$(h,k)$ integer selection for small $h + k$). A measurement of a capsid with T-number outside $\{1, 3, 4, 5, 7, 12, 13\}$ at small scales would challenge Proposition~\ref{prop:Tnumbers}.

%==================================================================
\section{Phyllotactic Divergence Angle}\label{sec:phyllotaxis}
%==================================================================

The phyllotactic divergence angle $\theta_{\mathrm{phyll}} \approx 137.508^\circ$ is observed in many plants (sunflower seed arrangement, pineapple spirals, etc.) and satisfies
\begin{equation}
\theta_{\mathrm{phyll}} = 360^\circ \cdot \bigl(1 - \varphi^{-1}\bigr) = \frac{360^\circ}{\varphi^2} = 137.5077640...^\circ.
\end{equation}

\begin{proposition}[Bio2: Phyllotactic angle from F1]\label{prop:phyllotaxis}
Under Paper XXXVI's F1 ($r = \varphi$, rigorous), the golden angle $360^\circ / \varphi^2$ is the unique ``maximally irrational'' rotation angle, i.e., the rotation angle at which sequential placements produce a dense, non-overlapping packing on a circle with slowest approach to any rational fraction. This is the mathematical reason why phyllotactic growth converges to this angle under an incremental-placement optimisation.
\end{proposition}

\begin{proof}[Argument]
Standard result of diophantine approximation~\citep{HardyWright}: $\varphi$ is the unique ``most irrational'' number in the sense of having the slowest-converging continued-fraction expansion ($\varphi = [1; 1, 1, 1, \ldots]$). Under Paper XXXVI F1, $\varphi$ is the cascade base permeability; hence the cascade-natural rotation is the golden angle, and phyllotactic growth under a local-packing objective converges to it.
\end{proof}

\begin{remark}[Independence of cascade]
Proposition~\ref{prop:phyllotaxis} is true whether or not one accepts the VFD cascade framework: it is a classical result about the golden ratio. The \emph{cascade-specific} content is that $\varphi$ is the structural permeability (F1), identifying the biological golden angle with the same $\varphi$ that sets mass hierarchies, $\Lambda$, and proton-radius predictions. The phyllotactic observation is therefore a structural consistency check: biology uses the same $\varphi$ the rest of the cascade does.
\end{remark}

%==================================================================
\section{DNA Helical Pitch}\label{sec:dna}
%==================================================================

B-form DNA has $\sim 10.5$ bp per turn (weak dependence on salt and sequence). The cascade-derivation notes (\citep{CascadeBioDNA}) propose a reading $10.5 \approx 2 \times \varphi^2 = 2 \times 2.618 = 5.236$ --- which does \emph{not} match. Alternative readings $10.5 = \varphi^5 - ?$ or $10.5 = \varphi^5 / 1.05$ are numerically close but structurally unclear.

\begin{remark}[DNA pitch remains an open target]
The programme currently does not have a clean cascade derivation of the DNA helical pitch. The observed $10.5$ bp/turn value has strong chemical-environmental dependence, and a purely geometric cascade derivation may be too crude. This is explicitly flagged as an open research question and not claimed as a cascade prediction in this paper.
\end{remark}

%==================================================================
\section{Falsifiable Predictions and Null Hypothesis}\label{sec:predictions}
%==================================================================

\begin{description}
\item[\textbf{B1}] Capsid T-numbers at small scales are drawn from $\{1, 3, 4, 5, 7, 12, 13\}$. \emph{Status}: common T-numbers fall in this set; rare large-T viruses ($T = 25, 27$) are outside and require deeper cascade interpretation.
\item[\textbf{B2}] Phyllotactic angle is $\varphi$-natural: $360^\circ / \varphi^2 \approx 137.508^\circ$. \emph{Status}: observed in many plants, independently of cascade framework.
\item[\textbf{B3}] DNA helical pitch: \emph{not} a cascade prediction in this paper.
\end{description}

\subsection{Null hypothesis}
Biological architecture is shaped by evolutionary pressure plus physical-chemical constraints; the cascade's role is to identify discrete constraints that natural selection cannot adjust. Where the cascade predicts a set (e.g., T-numbers), observed frequencies within the set should be consistent with combinatorial ubiquity; observed absences outside the set would challenge the cascade interpretation. Current data is consistent with the cascade as a \emph{constraint} rather than a \emph{driver} of biology.

%==================================================================
\section{Scope and Honest Limitations}\label{sec:scope}
%==================================================================

This is a skeleton paper. Its limits:
\begin{enumerate}
\item It does not derive any biological \emph{dynamics}. Only discrete architectural constraints are considered.
\item The T-number derivation is structural-proposal level, not rigorous.
\item DNA pitch is flagged as outside the paper's achievable scope.
\item No biological mechanism (e.g., capsid assembly, DNA binding, plant tropism) is modelled.
\end{enumerate}

\textbf{Dependencies}: Paper XXXVI F1 for the $\varphi$-scaling (phyllotaxis); Paper XXXV for the 600-cell $= 15 \times 40$ Hopf fibration.

\textbf{Future work}: a dedicated biology paper in a later programme phase would need to address (a) rigorous derivation of the T-number set from cascade structure; (b) DNA pitch from a chemical-geometric bridging mechanism; (c) quantitative null hypothesis testing against biological databases.

\subsection{Conclusion}

The cascade's biology rung places the $40$-cell icosahedral Hopf fibre as a structural locus for discrete biological architectures. Caspar--Klug T-numbers for small viruses are compatible with the cascade's $\twoI$-orbit selection. Phyllotactic golden angle is cascade-consistent via $\varphi$ (F1). DNA helical pitch remains an open target. The paper's contribution is honest: it maps the overlap between cascade geometry and biological architecture without overstating the explanatory reach.

%==================================================================
\bibliographystyle{plain}
\begin{thebibliography}{99}

\bibitem{SmartXXXV} L. Smart, \emph{QM, Life, and Gravity as $E_8$ Projections}, Paper XXXV, VFD Programme, 2026.
\bibitem{SmartXXXVI} L. Smart, \emph{Cascade Foundations: F1--F8}, Paper XXXVI, VFD Programme, 2026.
\bibitem{WOP43} L. Smart, \emph{WO-PAPER-XLIII --- Quantifying the Biology Rung}, VFD Programme work order, 2026.
\bibitem{CasparKlug1962} D. L. D. Caspar, A. Klug, \emph{Physical principles in the construction of regular viruses}, Cold Spring Harbor Symposia 27, 1--24 (1962).
\bibitem{HardyWright} G. H. Hardy, E. M. Wright, \emph{An Introduction to the Theory of Numbers}, 6th ed., Oxford University Press, 2008.
\bibitem{CascadeBioDNA} Cascade derivation note, \emph{cascade-bio.md} / \emph{cascade-genetic-code.md}, VFD Programme internal, 2024--2026.

\end{thebibliography}

\end{document}
